\section{Related Work}
\label{sec:related_work}

\subsection{Joint-Embedding Self-Supervised Learning}

Joint-embedding SSL methods learn representations by aligning features of augmented views of the same image. DINO \citep{caron2021emerging} and its successor DINOv2 \citep{oquab2023dinov2} use a student--teacher framework with cross-entropy over softmax probabilities; BYOL \citep{grill2020byol} uses a stop-gradient predictor; SimCLR \citep{chen2020simclr} uses a contrastive loss with negatives. All of these methods $\ell_2$-normalize the projection-head output, placing features on the unit hypersphere.

A parallel line of work enforces non-collapse via explicit regularization of the feature distribution rather than through the softmax-prototype or contrastive mechanisms. VICReg \citep{bardes2022vicreg} regularizes per-dimension variance and off-diagonal covariance; Barlow Twins \citep{zbontar2021barlow} decorrelates dimensions by penalizing off-diagonals of the cross-correlation matrix; MCR\textsuperscript{2} \citep{yu2020mcr2} maximizes the coding rate \citep{ma2007segmentation} of the representation. Most recently, SimDINO and SimDINOv2 \citep{wu2025simdino} remove the DINO head, softmax, centering, and sharpening entirely, replacing them with a squared-Euclidean alignment loss and a single coding-rate regularizer. SimDINO matches or exceeds DINO/DINOv2 performance with substantially less machinery. Our objective is a close variant of SimDINO adapted for unnormalized features: we use the coding rate on the correlation matrix \citep{kacker2025correlation_rate} rather than the covariance, and we add a two-sided variance penalty (SigReg \citep{schaeffer2024sigreg}) to pin the feature scale.

\subsection{Feature Norm as Confidence in SSL}

The closest prior work to ours is Draganov et al.~\citep{draganov2025embedding_norms}, who argue that in cosine-similarity SSL methods---specifically SimCLR, SimSiam, and BYOL---the unnormalized pre-projection embedding norm encodes a model-confidence signal. Their mechanism is density-based: frequently observed patterns receive more gradient updates and their embeddings grow larger, so typical samples end up with higher norms than atypical ones. They demonstrate the norm--accuracy coupling on several benchmarks and report that OOD inputs receive smaller norms. We share the headline finding in the norm-accuracy ranking, but our setting and scope differ. Their analysis covers \emph{spherical} SSL models whose training objective operates on $\ell_2$-normalized features; the norm they study is a by-product of the penultimate layer. We instead study \emph{flat-space} SSL, where the embedding norm enters the loss directly through the variance regularizer, and we characterize how the norm responds to continuous input corruption (\secref{sec:noise}) and where pure-noise inputs are placed in feature space (\secref{sec:noise_subspace})---neither of which is examined by Draganov et al. Our DINOv2 comparison (\secref{sec:dinov2}) also shows that the direction of the coupling is not automatic, even within spherical-SSL.

CSI \citep{tack2020csi} uses $\|\vect{z}\| \cdot \cos(\vect{z}, \vect{z}_\text{nn})$ as an OOD score on SimCLR features, exploiting the tendency of the contrastive objective to inflate in-distribution norms. CSI is a detection-time score built on top of an off-the-shelf contrastive model; our focus is the representation-level geometry itself, in a non-contrastive flat-space setting, without any labeled OOD data to calibrate against.

\subsection{Feature Norm as a Quality Signal}

In face recognition, several works have observed that the magnitude of a learned face embedding correlates with image quality. NormFace \citep{wang2017normface} studied the effect of $\ell_2$ normalization on softmax classifiers. MagFace \citep{meng2021magface} explicitly engineered a loss that makes the feature magnitude correlate with face image quality: low-quality faces (occlusion, blur) are placed near the origin, high-quality frontal faces far from it. AdaFace \citep{kim2022adaface} builds on this idea with a margin that adapts to the norm. Crucially, these works \emph{design for} the norm--quality coupling; they assume that if the loss is structured correctly, useful magnitude information will emerge. Our finding sits in the same conceptual space but arises without any explicit quality signal: the coupling emerges from the flat-Gaussian geometry alone.

\subsection{Energy-Based and Feature-Norm OOD Detection}

A separate literature uses norm-like quantities at test time to distinguish in-distribution from out-of-distribution inputs. Energy-based OOD \citep{liu2020energy} uses $-\logsumexp(\text{logits})$ as an OOD score, which reduces to a norm-like quantity when the classifier head is approximately linear. Mahalanobis-based OOD \citep{lee2018mahalanobis} uses class-conditional feature distributions. ViM \citep{wang2022vim} explicitly combines feature norm with softmax residual. SSD \citep{sehwag2021ssd} uses SimCLR features for OOD detection without labels. These methods typically assume a classifier has been trained and study the resulting feature geometry; our setting is purely self-supervised, and the OOD behavior is a consequence of the representation learning rather than an add-on.

\subsection{Representation Geometry and Dimensional Collapse}

Jing et al.~\citep{jing2022dimensional} study \emph{dimensional collapse} in contrastive SSL, where the learned features occupy a lower-dimensional subspace than the ambient space, usually characterized as a failure mode. Our observation that pure-noise features collapse to a low-dimensional subspace is closely related, but with a different interpretation: the collapse is \emph{selective}. In-distribution features span the full representation space; out-of-distribution inputs are the ones that collapse. This is the intended behavior of an OOD-aware representation, not a pathology.

Several works examine what is encoded in the radial and angular components of learned features \citep{wang2020hypersphere, wang2017normface}. Wang \& Isola \citep{wang2020hypersphere} decompose contrastive objectives into alignment (low angular spread within class) and uniformity (high angular spread across classes); our flat-space setting adds a radial degree of freedom that their framework cannot directly express.

\subsection{Summary of Position}

The literatures above---flat-space SSL, norm-as-confidence in cosine SSL \citep{draganov2025embedding_norms, tack2020csi}, face-recognition norm quality signals, OOD detection, and representation collapse---converge on our finding but have not previously been assembled for flat-space SSL. Draganov et al.~establish the norm-accuracy ranking for spherical SSL; we show the same ranking holds without $\ell_2$ normalization, and go further by characterizing the continuous noise-corruption response and the low-dimensional pure-noise subspace. The face-recognition community designs for norm--quality coupling; we show it emerges for free. The OOD community uses feature norm as a post-hoc detector; we show the detector is built in at representation-learning time. The SSL community has extensive work on avoiding representation collapse; we show that one kind of collapse (for OOD inputs) is a feature, not a bug, under the right geometry.
